\section{Conclusion}
\label{sec:conclusion}

We presented the first systematic comparison of information storage capacity in transformer hidden states across five encoding schemes spanning the full decoder-complexity spectrum. Our key finding is a clear \emph{decoder--capacity tradeoff}: raw bit-packing stores the most information (24,576 bits at fp32 in a 768-dimensional vector), but model-based decoders leverage learned structure to recover meaningful token sequences from far fewer bits. The unembedding matrix decodes random tokens perfectly from 3 dimensions per token but fails at 2, revealing a sharp geometric threshold. Superposition encoding, successful for sparse binary features, fails entirely for dense token IDs. A single frozen transformer layer provides almost no benefit beyond position embeddings, while the full 12-layer model dramatically improves decoding. Multi-vector scaling achieves 99.4\% accuracy on 500 random tokens with 32 vectors, with peak efficiency of 0.46 bits per learnable parameter at 16 vectors.

These results have practical implications for soft prompt design, memory token approaches, and model compression: the effective capacity of a hidden state depends as much on the decoder as on the vector itself. Future work should extend this analysis to larger models, diverse natural text, and hybrid encoding schemes that combine raw bit-packing with model-based decoding.
